

Modern DFT is built upon the foundation of the pivotal theorems established by Hohenberg and Kohn,\cite{kohn1999nobel} which states that the  Hamiltonian for a many-electron system is uniquely determined by its ground-state electron density $\rho(\bm{r})$, from which one can derive the corresponding system energy $E$. Most importantly, Hohenberg and Kohn demonstrated that the exact ground state energy can be determined by minimising the \textit{total energy functional} $E[\rho]$ over all possible electronic density, \emph{i.e.}
\begin{equation}\label{eq:DFT_min_rho}
E^{\mbox{\scriptsize{gs}}}=\min_{\rho}E[\rho(\bm{r})]    
\end{equation} 
This unique theory has provided the much needed simplification to tackle the many-body quantum mechanical problems for real materials by mapping an electronic wavefunction with $3n$ variables into the electron density $\rho(\bm{r})$ which is a function of only  three Cartesian coordinates. 

In order to make the Hohenberg-Kohn Theorems practical for performing realistic calculations, it is  necessary to write down the expressions for $E[\rho]$, and this is achieved later\cite{kohn1965self} with the introduction of the famous Kohn-Sham equation for DFT, in which the total energy functional is divided into the sum of multiple terms:
\begin{equation}\label{eq:KS_eqn}
E_{\mbox{\scriptsize{KS}}}[\rho]=\sum_{i=1}^{occ}\int\vert \nabla\phi^{\mbox{\scriptsize{KS}}}_i (\bm{r})\vert^2 \mbox{d}\bm{r} + \int v(\bm{r})\rho(\bm{r})\mbox{d}\bm{r}
+\frac{1}{2}\int \mbox{d}\bm{r} \int \mbox{d}\bm{r}^\prime \frac{\rho(\bm{r})\rho(\bm{r}^\prime)}{\left\vert \bm{r}-\bm{r}^\prime\right\vert} +E_{NN}+E_{xc}[\rho].
\end{equation}
Here, the first three terms (electron kinetic energy, electron-nuclear Coulomb interactions and electron-electron Coulomb interactions, respectively) can be written analytically as the functionals of the electronic density, that can be self-consistently determined by the occupied Kohn-Sham orbitals $\phi^{\mbox{\scriptsize{KS}}}(\bm{r})$, where $\rho(\bm{r})=\sum_{i=\mbox{\scriptsize{occ}}}|\phi^{\mbox{\scriptsize{KS}}}_i(\bm{r})|^2$. $E_{NN}$ is the Coulomb repulsions between the nuclei, which is deterministic from the nuclei charges and positions. These four terms, collectively, could be understood from the classical electrostatics, particularly based on the functional forms they adopt. In contrast, the \emph{exchange-correlation functional} $E_{xc}[\rho]$ describes the pure quantum mechanical nature of electrons, and can generally be separated into the sum of exchange and correlation contributions as $E_{xc}[\rho]=E_{x}[\rho]+E_{c}[\rho]$. 
Physically, the exchange energy is the energy associated with exchanging the positions of two electrons, while maintaining the antisymmetry of the electronic wavefunction according to the Pauli's exclusion principle. With this definition and utilising the Kohn-Sham orbitals $\phi^{\mbox{\scriptsize{KS}}}(\bm{r})$, the exchange energy, which is a \emph{non-local} quantity in $\bm{r}$, in principle, can be calculated exactly  using its definition from the Hartree-Fock (HF) theory:\cite{staroverov2012density}
\begin{equation}\label{eq:exchange_HF_DFT}
E_{x}^{\mbox{\scriptsize{exact}}}=-\sum_{i,j=1}^{n/2}\int \mbox{d}\bm{r}\int \mbox{d}\bm{r}^\prime  \frac{\phi_{i}(\bm{r}) \phi_{i}^{*}(\bm{r}^\prime) \phi_{j}^{*}(\bm{r})\phi_{j}(\bm{r}^\prime)}{\left\vert \bm{r}-\bm{r}^\prime\right\vert}.
\end{equation}
The negative sign in the above expression shows that the exchange interactions always contribute to the stabilisation of the system. On the other hand, electron correlations, which also contributes to stabilising forces and describes the extent to which the movement of one electron is affected by the presence and motion of all other electrons in the system, lack of a closed mathematical form such as \cref{eq:exchange_HF_DFT}. 

It turns out that simply appending approximated $E_{c}[\rho]$ to $E_{x}^{\mbox{\scriptsize{exact}}}$ do not generally lead to good performing results, because of the opposite effects on how the exchange and correlation interactions can affect the distributions of two particle electron density $\rho(\bm{r}_{1},\bm{r}_{2})$ in space. As such, \emph{the general lack of an analytical  form for  $E_{xc}[\rho]$, particularly one that is local, meaning only dependent on the electron density $\rho(\bm{r})$ at a single point in space}, is the major challenge behind DFT, and is feuling the ongoing development of DFAs to accurately approximate the true $E_{xc}$ for a many-electron system without significantly scarifying the computational efficacy of DFT. In what follows, we will briefly review  the major categories of DFAs that have been developed and matured over the years.